%% kapa: exam=1 round=37 year=2026 period=1 form=A subject=civil_law subject_ko=민법 seq=8 number=8
%% source: Q-Net 감정평가사 기출문제 (cst00309) artlSeq=5258924 file=raw/1008/5258924/2026년 제37회 감정평가사 1차 1교시.pdf
%% conversion: kordoc 4.24.51 -> tools/gen1.py (auto); review=proofread-sonnet; flags=none
\begin{question}{8}
물건에 관한 설명으로 옳은 것은? (다툼이 있으면 판례에 따름)
\begin{choices}
  \item 국립공원의 입장료는 민법상 과실이다.
  \item 물건의 용법에 의하여 수취하는 산출물은 법정과실이다.
  \item 부동산은 종물이 될 수 없다.
  \item 주물의 효용과 직접 관계가 없더라도 주물 소유자의 사용에 공여되는 물건은 종물이다.
  \item ‘종물은 주물의 처분에 따른다’는 법리는 압류와 같은 공법상 처분에도 적용된다.
\end{choices}
\end{question}
\answer{5}
